\documentclass[11pt]{article}
\usepackage[a4paper,margin=25mm,headheight=14pt]{geometry}
\usepackage[T1]{fontenc}
\usepackage{lmodern,microtype}
\usepackage{amsmath,amssymb,amsthm,mathtools}
\usepackage{booktabs,array,longtable}
\usepackage{xcolor,enumitem,fancyhdr}
\usepackage[colorlinks=true,linkcolor=blue!45!black,citecolor=blue!45!black,urlcolor=blue!45!black]{hyperref}
\usepackage{bookmark,xurl}
\newtheorem{theorem}{Theorem}[section]
\newtheorem{lemma}[theorem]{Lemma}
\newtheorem{corollary}[theorem]{Corollary}
\newtheorem{proposition}[theorem]{Proposition}
\theoremstyle{remark}\newtheorem{remark}[theorem]{Remark}
\newcommand{\R}{\mathbb R}
\newcommand{\norm}[1]{\left\lVert#1\right\rVert_2}
\newcommand{\abs}[1]{\left|#1\right|}
\newcommand{\fl}{\operatorname{fl}}
\newcommand{\PCG}{P_{\mathrm{CG}}}
\newcommand{\berr}{\operatorname{berr}}
\newcommand{\eps}{\varepsilon}
\newcommand{\diag}{\operatorname{diag}}
\newcommand{\spanof}{\operatorname{span}}
\newcommand{\coefnorm}[1]{\left\lVert#1\right\rVert_{\mathrm{coef},1}}
\newcommand{\energy}[1]{\left\lVert#1\right\rVert_A}
\pagestyle{fancy}
\fancyhf{}
\fancyhead[L]{\small IE-20: extended precision results}
\fancyhead[R]{\small Partial resolution}
\fancyfoot[C]{\thepage}
\setlength{\parindent}{0pt}
\setlength{\parskip}{4pt}
\setlist{itemsep=2pt,topsep=3pt}
\allowdisplaybreaks[2]
\hypersetup{pdftitle={IE-20: Exact Scalar Threshold and Extended Precision Bounds},pdfauthor={ChatGPT},pdfsubject={Extended partial resolution; the uniform matching-bound problem remains unresolved}}
\begin{document}
\hypersetup{pageanchor=false}
\begin{titlepage}
\vspace*{12mm}
{\large Numerical linear algebra\par}
\vspace{7mm}
{\LARGE\bfseries IE-20: Exact Scalar Threshold\\[3pt]and Extended Precision Bounds\par}
\vspace{7mm}
{\large Proofs, exact checks, and the remaining matching problem\par}
\vspace{10mm}
Prepared by ChatGPT\hfill September 13, 2026
\vspace{10mm}

\fbox{\begin{minipage}{0.93\textwidth}
\textbf{Completion status.} This is an \emph{extended partial resolution},
not a full solution to IE-20. It adds an exact scalar answer, a sharper
one-step result, a joint dimension/conditioning breakdown family,
a capped-condition-number upper bound, and two
finite-precision early-convergence arguments. The upper and lower bounds
still do not match for all parameter triples. Section~\ref{sec:gap}
gives a concrete remaining gap. The proofs are not independently refereed
or proof-assistant verified.
\end{minipage}}

\vspace{8mm}
\textbf{Abstract.}
We study the fixed CG recurrence and adversarial scalar relative-error
envelope specified by IE-20. The scalar threshold is determined exactly by
a degree-five rational inequality. A prefix-wise analysis sharpens the
dimension factor in the local matvec error. In general dimension, set
\[
H=\min\{K,64/\eps^2\},\qquad D=1024H^2,
\]
and let $m$ be the minimum of the dimension, an explicit Chebyshev horizon
of order $\sqrt K\log(K/\eps)$, and a backward-error horizon of order
$\eps^{-2/3}$. We prove the sufficient precision
\[
\left\lceil\log_2\frac{1024(n+1)K D^{12m+20}}{\eps^2}\right\rceil.
\]
The cap in $H$ follows from the fact that an excessively long step would
already have small backward error. The early-horizon arguments transfer
polynomial approximation to the actual rounded recurrence, rather than
assuming that computed iterates are exact Krylov projections.
A prefix-wise matvec estimate and a new breakdown family also give
\[
\PCG(n,K,\eps)=\log_2(nK)+O_{\eps}(1)
\]
for each fixed tolerance, uniformly in $n\ge2$ and $K\ge1$.
The additive constants are universal when $3/8\le\eps<1/2$.
There is also a sharp order uniformly for $1\le K\le2$. The previous path/outlier lower
bound and its exact certificates are preserved in the companion archive.

\vfill
\small
The package contains this PDF, editable \LaTeX, standard-library Python
code, regenerated exact checks, a proof-scope audit, source notes, and the
previous recovery archive. Finite checks verify their recorded examples,
not the universal analytic claims.
\end{titlepage}
\hypersetup{pageanchor=true}
\pagenumbering{roman}
\tableofcontents
\newpage
\pagenumbering{arabic}

\section{Problem, notation, and result map}
\label{sec:problem}

We use the editorial IE-20 formulation in the repository\cite{ie20}.
For $n\ge1$, $K\ge1$, and $0<\eps<1/2$, stored inputs are SPD
$A\in\R^{n\times n}$ and $b\ne0$, with $\kappa_2(A)\le K$ and at most
$p$ binary significant bits per entry. Exponents are unbounded. Each
scalar operation satisfies
\[
\fl(a\circ b)=(a\circ b)(1+\delta),\qquad |\delta|\le u=2^{-p}.
\]
Errors are independently admissible. Products and additions are separate;
dot products include all indices, accumulated from zero in increasing
order. Copies and comparisons are exact. The recurrence stops at a zero
recursive residual or before a zero divisor. Success requires a produced
iterate of index at most $n$ satisfying
\begin{equation}\label{eq:eta}
\eta(x)=\frac{\norm{b-Ax}}{\norm A\norm x+\norm b}\le\eps.
\end{equation}
The threshold $\PCG(n,K,\eps)$ must guarantee success at every greater
precision too, for every admissible input and error schedule.

\paragraph{Interpretation of the error model.}
The error equations are an envelope, not a correctly rounded binary map.
An addition of zero may carry an allowed relative error, and an exact
intermediate may be retained even if it is not a $p$-bit number. Input
representability is a separate condition. All lower-bound constructions
and the exact executor use this interpretation. A lower bound here is
not automatically a lower bound for IEEE round-to-nearest.

We write a search direction as $v_j$, reserving $p$ for precision:
\begin{align*}
&x_0=0,\quad r_0=v_0=b,\quad \rho_0=\fl(b^Tb),\\
&q_j=\fl(Av_j),\quad d_j=\fl(v_j^Tq_j),\quad
\alpha_j=\fl(\rho_j/d_j),\\
&x_{j+1}=\fl(x_j+\alpha_jv_j),\quad
r_{j+1}=\fl(r_j-\alpha_jq_j),\\
&\rho_{j+1}=\fl(r_{j+1}^Tr_{j+1}),\quad
\beta_j=\fl(\rho_{j+1}/\rho_j),\quad
v_{j+1}=\fl(r_{j+1}+\beta_jv_j).
\end{align*}
The last line is evaluated only when another step is needed. Every
product and addition inside a displayed vector expression is separately
rounded according to the preceding conventions.

\subsection{The strengthened upper bound}

Define the integer horizons
\begin{align}
\bar K&=\max\{2,K\},&
k_C&=\left\lceil\sqrt{\bar K}\right\rceil
      \left\lceil\log_2\frac{16\bar K}{\eps}\right\rceil,
\label{eq:kC}\\
k_U&=\left\lceil(64/\eps)^{2/3}\right\rceil,&
m&=\min\{n,k_C,k_U\},\label{eq:horizon}\\
H&=\min\{K,64/\eps^2\},&D&=1024H^2.\label{eq:HD}
\end{align}
The main upper-bound theorem, proved in Sections~\ref{sec:bootstrap}--\ref{sec:mainupper}, is
\begin{equation}\label{eq:mainupper}
\PCG(n,K,\eps)\le
\max\left\{2,\left\lceil\log_2
\frac{1024(n+1)K D^{12m+20}}{\eps^2}\right\rceil\right\}.
\end{equation}
Consequently, with universal constants,
\begin{equation}\label{eq:orderupper}
\PCG=O\!\left(
\log\frac{nK}{\eps}+
\min\!\left\{n,\sqrt K\log\frac K\eps,\eps^{-2/3}\right\}
\log\!\left(2\min\{K,\eps^{-2}\}\right)\right).
\end{equation}
The constants in \eqref{eq:mainupper} are conservative. The integer
horizon $k_C$ is deliberately larger than the least degree satisfying
the Chebyshev estimate.

\subsection{What is sharp, and what is not}

For $n=1$, Section~\ref{sec:scalar} gives the exact integer threshold.
For $n\ge2$, the general lower bound remains
\begin{equation}\label{eq:generallower}
\PCG(n,K,\eps)\ge c\left(\log_2\frac n\eps+\log_2K\right).
\end{equation}
Combining the new upper arguments with explicit witnesses yields
\begin{align}
\PCG(n,K,\eps)&=\Theta(\log_2(n/\eps)),&&1\le K\le2,
\label{eq:wellcond}\\
\PCG(n,K,\eps)&=\Theta(\log_2(nK)),&&n\ge2,\quad3/8\le\eps<1/2,
\label{eq:coarse}\\
\PCG(2,K,\eps)&=\Theta(\log_2K+\log_2(1/\eps)).
\label{eq:two}
\end{align}
All constants in these displayed $\Theta$ statements are universal on
the indicated domains. For each fixed tolerance, uniformly in $n\ge2$ and $K\ge1$, we prove
\begin{equation}\label{eq:fixedtol}
\PCG(n,K,\eps)=\log_2(nK)+O_{\eps}(1).
\end{equation}
In particular, for fixed $n\ge2$ and fixed $\eps$,
\begin{equation}\label{eq:largeK}
\PCG(n,K,\eps)=\log_2K+O_{n,\eps}(1).
\end{equation}
None of these statements closes the all-parameter gap between
\eqref{eq:generallower} and \eqref{eq:orderupper}.

\subsection{Relation to the recovered work and recent sources}

The prior manuscript in \path{prior/IE20_recovered.zip} supplied the
same elementary lower bounds, the $K=1$ and $n=2$ orders, an upper bound
with $H=K$, $m=n$, and an extra dimension factor in the local error budget, and a path/outlier family requiring order $n\log K$
on specified joint parameter sequences. This manuscript proves the new
upper statements directly, while retaining the old archive for the full
hard-family proof and provenance. No claim of external priority is made.

Derezi\'nski, Nakatsukasa, and Rebrova give an exact-arithmetic reciprocal
polynomial approximation used in Case 3 of Theorem~\ref{thm:main}\cite{minberr}.
The transfer to rounded CG iterates is proved here; it is not imported
as a finite-precision conclusion from their theorem.

Bake, Carson, and Ma establish error estimates at a sufficiently late
iteration under stated smallness conditions; their theorem does not bound
that iteration by $n$\cite{bcm}. A September 2026 manuscript by Sinha gives
a sufficient precision for a specified CG computation with independently
recomputed residual tests, together with a conditioning obstruction,
and explicitly retains a quantitative gap\cite{sinha}. These results do
not supply the missing matching characterization for the present target.

\section{An exact answer in dimension one}
\label{sec:scalar}

\begin{theorem}[Exact scalar threshold]\label{thm:scalar}
For every $K\ge1$, define
\begin{equation}\label{eq:scalarE}
E_1(u)=\frac{(1+u)^4-(1-u)^5}{(1+u)^4+(1-u)^5}.
\end{equation}
Then
\begin{equation}\label{eq:scalarthreshold}
\PCG(1,K,\eps)=\min\{p\in\mathbb Z:\ p\ge2,
\ E_1(2^{-p})\le\eps\}.
\end{equation}
In particular, the answer is independent of the supplied condition-number
bound $K$ and is computable by exact rational comparisons.
\end{theorem}
\begin{proof}
Write $A=[a]$, $a>0$, and $b\ne0$. Exactly nine elementary operations
contribute to the first iterate: two for $\rho_0$, two for $q_0$, two
for $d_0$, one for division, and two for the update of $x_1$.
All relevant divisors are positive because $u\le1/4$.
After cancellation of the input values,
\[
x_1=\frac ba F,
\]
where $F$ is a product of five allowed factors $1+\delta$ divided by
four such factors. Thus
\[
F_-:=\frac{(1-u)^5}{(1+u)^4}\le F\le
F_+:=\frac{(1+u)^5}{(1-u)^4}.
\]
For positive $F$, the backward error is $|1-F|/(1+F)$.
It increases with $|\log F|$, and
\[
-\log F_- -\log F_+=-\log(1-u^2)>0.
\]
The worse endpoint is therefore $F_-$. Its error is precisely
\eqref{eq:scalarE}. This endpoint is attained with the admissible inputs
$a=b=1$, errors $-u$ in the five numerator operations and $+u$ in the
four denominator operations. The initial iterate has error one.
Residual-update operations occur after $x_1$ has been produced and cannot
alter its backward error.

The function $F_-(u)$ is strictly decreasing on $(0,1/4]$, so $E_1(u)$
is strictly increasing. Consequently the least successful precision also
guarantees success at every larger precision. This proves the threshold
quantifier in \eqref{eq:scalarthreshold}.
\end{proof}

The leading small-$u$ behavior is $E_1(u)=(9/2)u+O(u^2)$, so
$\PCG(1,K,\eps)=\log_2(1/\eps)+O(1)$. The exact formula, rather than
this expansion, handles equality cases and integer rounding. The
companion code exhausts all $2^9$ sign corners at each of three precisions;
the analytic monotonicity argument covers the entire continuous error box.

\section{Elementary lower bounds retained for comparison}
\label{sec:lower}

\begin{theorem}\label{thm:lower}
For every admissible triple,
\begin{equation}\label{eq:loweridentity}
\PCG(n,K,\eps)\ge
\max\left\{2,1+\left\lfloor\log_2\frac{n}{4\eps}\right\rfloor\right\}.
\end{equation}
For $n\ge2$ there is also the bound
\begin{equation}\label{eq:lowerbreakdown}
\PCG(n,K,\eps)\ge\max\{2,\lfloor\log_2(K+1)\rfloor\}.
\end{equation}
Together they imply \eqref{eq:generallower}.
\end{theorem}
\begin{proof}
For \eqref{eq:loweridentity}, choose $A=I_n$ and $b=e_1$.
Set the error to $+u$ on every addition in the first row of the first
matrix-vector product, including additions of zero, and set every other
error to zero. Then $q_0=t e_1$ with $t=(1+u)^n$. The produced iterate
and recursive residual are $x_1=t^{-1}e_1$ and $r_1=0$. The algorithm
stops, but its true backward error is $(t-1)/(t+1)$.
For $u\le1/4$ and $\eps<1/2$,
\[
\log(1+u)\ge\tfrac45u,\qquad
\log\frac{1+\eps}{1-\eps}\le\tfrac83\eps.
\]
Hence $u\ge4\eps/n$ gives strict failure. Failure at
$p=\lfloor\log_2(n/(4\eps))\rfloor\ge2$ forces the threshold to exceed
$p$. Smaller values are covered by its defining lower bound two.

For \eqref{eq:lowerbreakdown}, at precision $p\ge2$ set $u=2^{-p}$ and
\[
A=\begin{pmatrix}1&1-u\\1-u&1\end{pmatrix}\oplus I_{n-2},\qquad
b=(1,-1,0,\ldots,0)^T.
\]
The stored number $1-u$ has exactly $p$ significant bits. The condition
number is $(2-u)/u=2^{p+1}-1$. Put error $-u$ on the two diagonal
products in the first matrix-vector product and zero elsewhere. Both
nonzero rows now sum to zero, so $q_0=d_0=0$. Only $x_0$ is produced,
and it fails. If $q=\lfloor\log_2(K+1)\rfloor\ge3$, take $p=q-1$;
then this input has condition number at most $K$ and the threshold is at
least $q$. The remaining cases again follow from $\PCG\ge2$.

Finally, the maximum of the two logarithmic lower bounds is at least
half their sum, up to additive constants absorbed by $\PCG\ge2$.
\end{proof}


\begin{theorem}[Joint dimension/conditioning breakdown]\label{thm:jointlower}
Let $n\ge2$, $K\ge15$, and $a=\lfloor n/2\rfloor$. Then
\begin{equation}\label{eq:jointinteger}
\PCG(n,K,\eps)\ge
1+\left\lfloor\log_2\frac{a(K+1)}2\right\rfloor.
\end{equation}
For all $n\ge2$, $K\ge1$, and admissible tolerances, this implies
\begin{equation}\label{eq:jointlower}
\PCG(n,K,\eps)\ge\log_2(nK)-5.
\end{equation}
\end{theorem}
\begin{proof}
Set $p=\lfloor\log_2(a(K+1)/2)\rfloor$, $u=2^{-p}$,
$t=au$, and $c=1-t$. These choices give
\[
p\ge2,\qquad \frac2{K+1}\le t<\frac4{K+1}\le\frac14,
\qquad nu<1.
\]
For the last estimate, use $u<4/[a(K+1)]$ and $a\ge n/4$.
Put the block $\bigl[\begin{smallmatrix}1&c\\c&1\end{smallmatrix}\bigr]$
on indices $1$ and $n$, put identity on the remaining indices, and
choose $b=e_1-e_n$. The condition number is $(2-t)/t\le K$.
The stored entry $c=(2^p-a)/2^p$ has at most $p$ significant bits.

In each active row, the first nonzero product is followed by $n-1$
additions before the final column is included. Thus one product and
$n-1$ additions supply $n$ independently selectable factors. In row 1,
scale its initial value one down to $c$; in row $n$, scale its initial
value $c$ up to one. Both scalings are admissible. Indeed,
\[
t=au\le\frac{nu}{2}\le\frac{nu}{1+nu},\qquad
c(1+u)^n\ge c(1+nu)\ge1.
\]
Thus $c\ge(1+u)^{-n}\ge(1-u)^n$ as well.
For an exact rational schedule, repeatedly use factor $1-u$ on the
first row until the next such factor would undershoot $c$, then use the
exact factor needed to reach $c$ and subsequently factor one. On the
last row use the corresponding greedy factors up to $1+u$ toward one.
Every chosen factor is rational and lies in the permitted interval.

Evaluate the last-column products and additions exactly. They cancel
$c-c$ in row 1 and $1-1$ in row $n$. All other rows are zero, so
$q_0=d_0=0$. The run breaks down with only $x_0$ produced. Failure at
precision $p$ proves \eqref{eq:jointinteger}.
Since $a\ge n/4$, this lower bound exceeds $\log_2(n(K+1)/8)$.
For $K<15$, the identity witness and $\eps<1/2$ give
$\PCG>\log_2(n/2)>\log_2(nK)-5$. Combining the two ranges proves
\eqref{eq:jointlower}. Together with the identity witness this also gives
\[
\PCG(n,K,\eps)\ge
\log_2\left(n\max\{K,\eps^{-1}\}\right)-5.
\]
\end{proof}

\section{The sharp exact first-step envelope}
\label{sec:first}

\begin{theorem}\label{thm:firstexact}
For exact arithmetic on real SPD inputs with $n\ge2$ and condition number
at most $K$, the largest possible first-step backward error is
\begin{equation}\label{eq:M}
M(K)=\frac{K-1}{\sqrt{8K(K+1)}}.
\end{equation}
This is a real-input extremal statement; exact representability of its
maximizing vector at a prescribed finite precision is not asserted.
\end{theorem}
\begin{proof}
Normalize $\norm A=\norm b=1$, put $a=1/K$, and write
$\mu=b^TAb$, $\nu=b^TA^2b$. The exact first iterate is $x=b/\mu$, so
\[
\eta(x)^2=\frac{\nu-\mu^2}{(1+\mu)^2}.
\]
Every eigenvalue $t\in[a,1]$ satisfies
$t^2\le(1+a)t-a$. Consequently
$\nu-\mu^2\le(1-\mu)(\mu-a)$. The identity
\[
\frac{(1-a)^2}{8(1+a)}-
\frac{(1-t)(t-a)}{(1+t)^2}
=
\frac{\bigl((a+3)t-(3a+1)\bigr)^2}
     {8(1+a)(1+t)^2}
\]
shows that the maximum occurs at $\mu=(1+3a)/(3+a)$.
The resulting square root is \eqref{eq:M}.
For $K>1$, equality is attained by a two-point spectrum $\{a,1\}$ and
a unit vector whose squared weights are $2/(3+a)$ at $a$ and
$(1+a)/(3+a)$ at one. Extra coordinates may have zero right-hand-side
weight. For $K=1$, the error is identically zero.
\end{proof}

\begin{theorem}[Finite-precision first step]\label{thm:firstfinite}
If $\eps>M(K)$, then
\begin{equation}\label{eq:firstprecision}
\PCG(n,K,\eps)\le
\max\left\{2,\left\lceil\log_2
\frac{512(n+1)K}{\eps-M(K)}\right\rceil\right\}.
\end{equation}
The conclusion also holds for $n=1$, although Theorem~\ref{thm:scalar}
is sharper there.
\end{theorem}
\begin{proof}
Continue with $\norm A=\norm b=1$. Set
$\omega=64(n+1)Ku$ and suppose $\omega\le1/16$.
The local scalar estimates proved in Lemma~\ref{lem:local} give
$\alpha_0\mu=1+\theta$, $|\theta|\le\omega$, and ensure a positive
computed denominator. The separately rounded product and addition in
$x_1$ then give
\[
\norm{x_1-b/\mu}\le\frac{2\omega}{\mu}.
\]
Indeed their additional relative deviation is at most $2u+u^2$, and
$u\le\omega/128$. Perturbing the numerator and denominator of
\eqref{eq:eta}, with $x^{\rm ex}=b/\mu$, yields
\[
\eta(x_1)\le\frac{\eta(x^{\rm ex})+2\omega}{1-2\omega}
\le M(K)+4\omega.
\]
For the last inequality, use $M(K)<1/(2\sqrt2)$ and $2\omega\le1/8$.
Let $\Delta=\eps-M(K)>0$. The precision in \eqref{eq:firstprecision}
implies $\omega\le\Delta/8<1/16$, so
$\eta(x_1)\le M(K)+\Delta/2<\eps$. This holds at all larger precisions
as well.
\end{proof}

\begin{corollary}\label{cor:first}
For $n\ge2$ and $3/8\le\eps<1/2$, \eqref{eq:coarse} holds.
Also, whenever $n\ge2$ and $\eps\ge2M(K)$,
\[
\PCG(n,K,\eps)=\Theta\!\left(\log_2\frac{nK}{\eps}\right)
\]
with universal constants on this region.
\end{corollary}
\begin{proof}
Since $1/(2\sqrt2)<23/64$, the coarse-tolerance region has
$\eps-M(K)\ge1/64$. Equation~\eqref{eq:firstprecision} is then at most
$\lceil\log_2(32768(n+1)K)\rceil$.
The joint lower bound below improves this to the explicit additive comparison
\[
\log_2(nK)-5\le\PCG(n,K,\eps)\le\log_2(nK)+17.
\]
In the second region, $\eps-M(K)\ge\eps/2$, and the same comparison
proves the stated order.
\end{proof}

\section{A failure-conditioned bootstrap}
\label{sec:bootstrap}

This section supplies the common recurrence analysis for all three
horizons in \eqref{eq:horizon}. The important new point is that $H$ in
\eqref{eq:HD}, rather than $K$, controls the step lengths while every
produced iterate fails.

\subsection{Normalization and local error equations}

Homogeneity allows $\norm A=\norm b=1$. To see this without assuming
representability of normalized inputs, let $a_0=\norm A$ and
$B_0=\norm b$ and transform an execution by
\[
A'=A/a_0,\quad b'=b/B_0,\quad x_j'=a_0x_j/B_0,\quad
r_j'=r_j/B_0,\quad v_j'=v_j/B_0.
\]
Then $\alpha_j'=a_0\alpha_j$, $\rho_j'=\rho_j/B_0^2$, and $\beta_j$
is unchanged. Every scalar error equation transforms with the same
$\delta$. The normalized analysis applies to all real inputs, hence also
to this possibly nonrepresentable normalized input. The backward error
is invariant. The normalized spectrum lies in $[1/K,1]$.

Set
\begin{equation}\label{eq:omega}
\omega=64(n+1)Ku,
\end{equation}
and write
\[
R_j=\norm{r_j},\quad V_j=\norm{v_j},\quad X_j=\norm{x_j},\quad
Q_j=v_j^TAv_j.
\]
The symbol $Q_j$ here is a scalar quadratic form, not a Lanczos basis.


\begin{lemma}[Normwise prefix estimate]\label{lem:prefix}
For the specified left-to-right matvec, if $nu\le1/2$, then
\begin{equation}\label{eq:prefixbound}
\norm{\fl(Av)-Av}\le4nu\norm A\norm v.
\end{equation}
The estimate does not replace $A$ by its entrywise absolute value.
\end{lemma}
\begin{proof}
Collect the row accumulators at their common column index $j$ into a
vector $s_j$. This is an algebraic grouping of a completed execution;
it does not change the row-by-row evaluation order or its errors.
Let $A_{:,j}$ be column $j$ and let
$y_j=\sum_{i=1}^j A_{:,i}v_i$ be its exact prefix. Diagonal matrices
$D_j^{\rm add},D_j^{\rm mul}$ of norm at most $u$ encode the componentwise
errors. For $e_j=s_j-y_j$,
\[
e_j=D_j^{\rm add}y_j+
(I+D_j^{\rm add})(e_{j-1}+D_j^{\rm mul}A_{:,j}v_j).
\]
Unrolling from $e_0=0$ gives
\[
\norm{e_n}\le(1+u)^n u\left(
\sum_{j=1}^n\norm{y_j}+\sum_{j=1}^n\norm{A_{:,j}}\,|v_j|\right).
\]
Each exact prefix has norm at most $\norm A\norm v$, while the second
sum is at most $\lVert A\rVert_F\norm v\le\sqrt n\norm A\norm v$.
Since $(1+u)^n\le e^{nu}<2$, the result follows from
$2(n+\sqrt n)\le4n$. All zero terms remain included.
\end{proof}

\begin{lemma}[Local envelopes]\label{lem:local}
Suppose $\omega\le1/2$. Whenever the relevant nonzero residuals and
directions are available, their computed scalar divisors are positive,
and the performed operations satisfy
\begin{align}
\alpha_jQ_j&=R_j^2(1+\theta_j),&|\theta_j|&\le\omega,\label{eq:localalpha}\\
\beta_jR_j^2&=R_{j+1}^2(1+\zeta_j),&|\zeta_j|&\le\omega,\label{eq:localbeta}\\
r_{j+1}&=r_j-\alpha_jAv_j+f_j,&
\norm{f_j}&\le\omega(R_j+\alpha_jV_j),\label{eq:localr}\\
x_{j+1}&=x_j+\alpha_jv_j+g_j,&
\norm{g_j}&\le\omega(X_j+\alpha_jV_j),\label{eq:localx}\\
v_{j+1}&=r_{j+1}+\beta_jv_j+h_j,&
\norm{h_j}&\le\omega(R_{j+1}+\beta_jV_j).\label{eq:localv}
\end{align}
Equations involving the next direction are used only when it is formed.
\end{lemma}
\begin{proof}
Let $\gamma=(n+1)u/[1-(n+1)u]$. Expanding at most $n+1$ factors
$1+\delta$ in each dot-product term gives
$\gamma\le2(n+1)u$. The assumption $\omega\le1/2$ implies
$(n+1)Ku\le1/128$, so Lemma~\ref{lem:prefix} gives
\[
\norm{q_j-Av_j}\le4nuV_j.
\]
The subsequent dot product satisfies
\begin{align*}
|d_j-Q_j|
&\le4nuV_j^2+2(n+1)uV_j\norm{q_j}\\
&\le8(n+1)uV_j^2\le8(n+1)KuQ_j.
\end{align*}
Its relative error is at most $\omega/8\le1/16$, so $d_j>0$ for
$v_j\ne0$. Also $|\rho_j-R_j^2|\le\gamma R_j^2$, with relative error
at most $\omega/32$, and $\rho_j>0$ for $r_j\ne0$.
The rounded division has error at most $u\le\omega/128$.
Taking the two ratios therefore proves \eqref{eq:localalpha} and
\eqref{eq:localbeta}, with bounds strictly smaller than $\omega$.
A separately rounded product and vector addition contribute at most
$3u$ times the sum of the input norms. Combining this observation with
the matvec estimate proves \eqref{eq:localr}--\eqref{eq:localv}.
All estimates retain the additions of zero required by the specification.
\end{proof}

\subsection{Induction statement}

Fix an integer $1\le m\le n$, and use $H,D$ from \eqref{eq:HD}.
Assume
\begin{equation}\label{eq:smallness}
\sigma^2:=\omega D^{12m+20}\le\eps^2/16.
\end{equation}
Let $E_j$ be the maximum absolute off-diagonal inner product among
$r_0,\ldots,r_j$, and the maximum absolute off-diagonal $A$-inner
product among formed directions through $v_j$. At the endpoint $j=m$,
include only directions through $v_{m-1}$; $v_m$ is not needed.
Empty maxima are zero. Put $B=32H$ and
\[
w_j=b-Ax_j-r_j,\qquad W_j=\norm{w_j}.
\]

\begin{lemma}[Bootstrap under total failure]\label{lem:bootstrap}
Assume that every iterate produced through the horizon fails
\eqref{eq:eta}. Under \eqref{eq:smallness}, the execution cannot stop
before producing $x_m$. For $0\le j\le m$, it satisfies
\begin{align}
\max\{1,R_j,X_j,V_j\}&\le D^{j+1},\label{eq:Mbound}\\
E_j&\le\omega D^{8j+8},\label{eq:Ebound}\\
W_j&\le4\omega D^{j+2}\le\eps/8,\label{eq:Wbound}\\
R_j&\ge\eps X_j+7\eps/8.\label{eq:failR}
\end{align}
The terms involving $V_j$ concern only $j<m$. Every needed coefficient
obeys
\begin{equation}\label{eq:abbounds}
\frac14\le\alpha_j\le8H,\qquad0\le\beta_j\le B,
\end{equation}
and the corresponding local errors obey
\begin{equation}\label{eq:errbounds}
\norm{f_j},\norm{g_j}\le\omega D^{j+2},\qquad
\norm{h_j}\le2\omega D^{j+2}.
\end{equation}
\end{lemma}

\begin{proof}
The initial conditions have $R_0=V_0=1$, $X_0=W_0=E_0=0$.
At each subsequent index, \eqref{eq:Wbound} and failure imply
\[
R_j\ge\norm{b-Ax_j}-W_j>
\eps(X_j+1)-W_j\ge\eps X_j+7\eps/8.
\]
Thus $R_j>\sigma$. We establish the remaining estimates in the order
below; this order is important for excluding breakdown without assuming
that all steps already exist.

\paragraph{The current direction is nonzero.}
Expanding \eqref{eq:localv} gives a combination of residuals through
$r_j$ with coefficient sum at most $2B^j$, plus accumulated local errors.
Since $B/D\le1/32$, their total norm is at most
$3\omega D^{j+1}$. Consequently
\begin{align*}
|r_j^Tv_j-R_j^2|
&\le2D^jE_j+3\omega D^{2j+2}\\
&\le\omega D^{9j+10}\le\sigma^2/8.
\end{align*}
For $j<m$, the last comparison follows from \eqref{eq:smallness} and
$D\ge1024$. Since $R_j>\sigma$, it gives $V_j\ge7R_j/8>0$.
Lemma~\ref{lem:local} now ensures $d_j>0$, so the next step can legally
be evaluated.

\paragraph{The upper bound on the step length.}
If $H=K$, use $Q_j\ge V_j^2/K$ in \eqref{eq:localalpha} to obtain
$\alpha_j\le8H$.
Suppose instead that $H=64/\eps^2$. Let $\alpha=\alpha_j$,
$s=\alpha v_j$, and $S=\norm s$. The induction gives
$X_j\le R_j/\eps$, $W_j\le R_j/7$, and $S\ge7\alpha R_j/8$.
If $\alpha>8H=512/\eps^2$, these inequalities imply $X_j/S\le1/4$.
Also $A^2\preceq A$ and \eqref{eq:localalpha} give
\[
\norm{As}\le\sqrt{2\alpha}\,R_j,
\qquad \norm{g_j}\le\omega(X_j+S).
\]
The denominator in the next backward error is at least $S/2$.
Using $b-Ax_j=r_j+w_j$ in its numerator therefore yields the safe bound
\begin{equation}\label{eq:largestepsuccess}
\eta(x_{j+1})\le\frac5\alpha+\frac6{\sqrt\alpha}+3\omega<\eps.
\end{equation}
The last inequality follows from $\alpha>512/\eps^2$,
$\eps<1/2$, and $\omega\le\eps/32$, all consequences of our assumptions.
This contradicts failure of the newly produced iterate. Hence
$\alpha_j\le8H$ in the capped case too. The argument has not presumed
bounds on $f_j,g_j$ beyond the local envelopes, nor presumed existence
of a later coefficient.

\paragraph{The lower bound on the step length.}
For $j\ge1$, write $v_j=r_j+\beta_{j-1}v_{j-1}+h_{j-1}$ and expand the
quadratic form in the opposite direction. With
$C=v_{j-1}^TAv_j$, this gives
\begin{align*}
Q_j={}&r_j^TAr_j+2\beta_{j-1}C-
\beta_{j-1}^2v_{j-1}^TAv_{j-1}\\
&+2h_{j-1}^TA(v_j-\beta_{j-1}v_{j-1})-h_{j-1}^TAh_{j-1}\\
\le{}&R_j^2+2BE_j+
2\norm{h_{j-1}}(V_j+B V_{j-1})\\
\le{}&R_j^2+\omega D^{8j+10}
\le R_j^2+\sigma^2/8\le2R_j^2.
\end{align*}
For $j=0$, $Q_0\le R_0^2$ directly. Equation~\eqref{eq:localalpha}
therefore gives $\alpha_j\ge1/4$.

\paragraph{Next residual and iterate.}
The two update envelopes imply
\[
R_{j+1},X_{j+1}\le(1+\omega)(1+8H)D^{j+1}
\le\tfrac14D^{j+2},\qquad
\norm{f_j},\norm{g_j}\le\omega D^{j+2}.
\]
Their true-residual gap satisfies $w_{j+1}=w_j-Ag_j-f_j$.
Summing gives \eqref{eq:Wbound} at the next index. If the new recursive
residual were zero, this gap bound would already contradict failure.
Thus a needed next norm-square divisor is positive.

\paragraph{Bounding the next ratio.}
Using $A^2\preceq A$ and
$r_j^TAv_j=Q_j-\beta_{j-1}v_{j-1}^TAv_j-h_{j-1}^TAv_j$ in the
squared residual update gives
\begin{align*}
R_{j+1}^2\le{}&R_j^2+(\alpha_j-2)R_j^2(1+\theta_j)
+2\alpha_jB E_j+2\alpha_j\norm{h_{j-1}}V_j\\
&+2\norm{f_j}(R_j+\alpha_jV_j)+\norm{f_j}^2\\
\le{}&12H R_j^2+D E_j+\omega D^{2j+4}\\
\le{}&12H R_j^2+\omega D^{8j+10}\le13H R_j^2.
\end{align*}
At $j=0$ the terms with index $-1$ are absent. For explicit constant
checks, $2\alpha_jB\le D/2$, while the other local terms are bounded by
$235H^2\omega D^{2j+2}$, which is less than
$\omega D^{2j+4}$. Equation~\eqref{eq:localbeta} yields
$\beta_j\le32H=B$. Then
\[
V_{j+1}\le(1+\omega)(R_{j+1}+B V_j)
\le\tfrac32(\tfrac14+\tfrac1{32})D^{j+2}<D^{j+2}.
\]
The bound on $h_j$ follows immediately.

\paragraph{Residual orthogonality defects.}
Put
$F_j=\max\{E_j,|r_i^Tr_{j+1}|:0\le i\le j\}$.
For $i<j$, substitute
$r_i=v_i-\beta_{i-1}v_{i-1}-h_{i-1}$ into
$r_i^Tr_{j+1}=r_i^Tr_j-\alpha_jr_i^TAv_j+r_i^Tf_j$.
For $i=j$, use
\[
r_j^Tr_{j+1}=R_j^2-\alpha_jQ_j+
\alpha_j\beta_{j-1}v_{j-1}^TAv_j+
\alpha_jh_{j-1}^TAv_j+r_j^Tf_j.
\]
In each case the coefficient of $E_j$ is below $D$, and the remaining
terms are bounded by $\omega D^{2j+5}$. Hence
\begin{equation}\label{eq:Fbound}
F_j\le D E_j+\omega D^{2j+5}.
\end{equation}

\paragraph{Direction conjugacy defects.}
When $v_{j+1}$ is needed, use
$Av_i=(r_i-r_{i+1}+f_i)/\alpha_i$. For $i<j$, this gives
\[
|v_i^TAv_{j+1}|\le8F_j+B E_j+2\omega D^{2j+3}.
\]
For the adjacent pair the exact algebraic cancellation is
\begin{align*}
\alpha_jv_j^TAv_{j+1}
={}&r_j^Tr_{j+1}-R_{j+1}^2+f_j^Tr_{j+1}
+\alpha_j\beta_jQ_j+\alpha_jv_j^TAh_j,\\
\alpha_j\beta_jQ_j
={}&R_{j+1}^2(1+\theta_j)(1+\zeta_j).
\end{align*}
Therefore
$|v_j^TAv_{j+1}|\le4F_j+2\omega D^{2j+4}$.
Combining with \eqref{eq:Fbound} proves
\[
E_{j+1}\le D^2E_j+\omega D^{2j+7}
\le\omega\bigl(D^{8j+10}+D^{2j+7}\bigr)
\le\omega D^{8(j+1)+8}.
\]
At the horizon no next direction is required, so \eqref{eq:Fbound}
alone suffices. This completes the induction and the exclusion of every
premature stopping mechanism under the failure assumption.
\end{proof}

\section{Transfer from polynomials to computed spans}
\label{sec:span}

For a polynomial $\pi(t)=\sum c_it^i$, let
$\coefnorm{\pi}=\sum|c_i|$. The next lemma is an explicit substitute for
an unjustified claim that rounded CG spans an exact Krylov space.

\begin{lemma}[Approximate polynomial containment]\label{lem:span}
Under Lemma~\ref{lem:bootstrap}, fix $1\le k\le m$. If $\deg\pi<k$ and
$\coefnorm{\pi}\le D^{k+1}$, then $y=\pi(A)b$ has a representation
\begin{equation}\label{eq:spanv}
y=\sum_{i=0}^{k-1}c_i v_i+e,\qquad
\sum|c_i|\le D^{2k},\qquad\norm e\le\omega D^{3k+1}.
\end{equation}
There is also a vector $z\in\spanof\{x_1,\ldots,x_k\}$ such that
\begin{equation}\label{eq:spanx}
\norm{z-y}\le\omega D^{3k+2}.
\end{equation}
Furthermore, for $i<k$,
\begin{equation}\label{eq:rvorth}
|r_k^Tv_i|\le\omega D^{9k+10}.
\end{equation}
\end{lemma}
\begin{proof}
For $i\le k-2$, the local equations give
\begin{equation}\label{eq:Avspan}
Av_i=\frac{(1+\beta_i)v_i-\beta_{i-1}v_{i-1}-v_{i+1}}{\alpha_i}
+\frac{f_i+h_i-h_{i-1}}{\alpha_i},
\end{equation}
with negative-index terms zero. The coefficient sum is at most
$8+256H\le D$, and the error norm is at most $\omega D^{i+3}$.
Inductively,
\[
A^t b=\sum_{i=0}^t c_{ti}v_i+e_t,\qquad
\sum_i|c_{ti}|\le D^t,\qquad
\norm{e_t}\le2\omega D^{2t+1}.
\]
The error induction uses $\norm A=1$:
$\norm{e_{t+1}}\le2\omega D^{2t+1}+\omega D^{2t+3}
\le2\omega D^{2t+3}$.
Combining the monomials of $\pi$ yields \eqref{eq:spanv}.

Next, substitute
$v_i=(x_{i+1}-x_i-g_i)/\alpha_i$ in the direction representation.
Since $x_0=0$, the error-free part is in the claimed iterate span. Its
additional error is at most
$4D^{2k}\omega D^{k+1}$; together with $e$ this is below
$\omega D^{3k+2}$.

Finally, expand $v_i$ into earlier residuals as in the bootstrap proof.
The residual coefficients have sum at most $2D^i$ and the accumulated
local error has norm at most $3\omega D^{i+1}$. Since $k>i$,
\[
|r_k^Tv_i|\le2D^i E_k+3\omega D^{i+1}R_k
\le\omega D^{9k+10}.
\]
\end{proof}

\section{The three termination arguments}
\label{sec:mainupper}

\begin{theorem}[Capped-condition, early-horizon upper bound]\label{thm:main}
With the definitions \eqref{eq:kC}--\eqref{eq:HD}, every precision
satisfying \eqref{eq:smallness} guarantees a successful produced iterate
by step $m$. Therefore \eqref{eq:mainupper} and \eqref{eq:orderupper} hold.
\end{theorem}

\begin{proof}
Assume for contradiction that all produced iterates fail. By
Lemma~\ref{lem:bootstrap}, every iterate through $m$ exists and the
bootstrap estimates hold. At least one of the following three cases
applies to the minimum defining $m$.

\subsection*{Case 1: the dimension horizon $m=n$}
Normalize $r_0,\ldots,r_n$ to unit vectors. Their Gram matrix has diagonal
one and off-diagonal entries of magnitude at most
\[
\frac{E_n}{\sigma^2}\le D^{-4n-12}.
\]
Every off-diagonal row sum is smaller than one. The inequality
$2|a_i a_j|\le a_i^2+a_j^2$ shows directly that this Gram matrix is
positive definite. Its rank would be $n+1$, impossible for vectors in
$\R^n$.

\subsection*{Case 2: the Chebyshev horizon $m=k_C$}
Put $k=m$, $a=1/\bar K$, and
\[
q(t)=\frac{T_k((1+a-2t)/(1-a))}{T_k((1+a)/(1-a))},\qquad
\pi(t)=\frac{1-q(t)}t.
\]
Since $q(0)=1$, $\pi$ is a polynomial of degree less than $k$.
The affine argument of $T_k$ has coefficient sum at most seven because
$a\le1/2$. The Chebyshev recurrence gives
\[
\coefnorm{T_k((1+a-2t)/(1-a))}\le15^k.
\]
Its normalizing denominator is at least one. Thus
\[
\coefnorm\pi\le16^k\le D^{k+1}.
\]
On the spectrum of $A$,
\[
|q(t)|\le2\vartheta^k,\qquad
\vartheta=\frac{\sqrt{\bar K}-1}{\sqrt{\bar K}+1}.
\]
For $x_*=A^{-1}b$ and $y=\pi(A)b$, it follows that
\begin{equation}\label{eq:exactenergy}
\energy{x_*-y}^2\le4K\vartheta^{2k}.
\end{equation}

We now compare $y$ with the actual rounded $x_k$, not an exact CG
projection. Let $\bar x_k=\sum_{i<k}\alpha_i v_i$, so
$\norm{x_k-\bar x_k}\le2\omega D^{k+1}$.
By Lemma~\ref{lem:span}, the difference between the direction
representation of $y$ and $\bar x_k$ has coefficient sum at most
$D^{2k+1}$, and the remaining vector error is at most
$\omega D^{3k+2}$. Combining \eqref{eq:rvorth} with the residual gap
therefore gives
\begin{align}
|(b-Ax_k)^T(y-x_k)|
&\le\omega D^{11k+11}+\omega D^{4k+3}
     +8\omega D^{2k+3}\notag\\
&\le\omega D^{11k+12}.\label{eq:crossbound}
\end{align}
For the last gap term, we used
$\norm y\le\coefnorm\pi\le D^{k+1}$ and $X_k\le D^{k+1}$;
this bound does not replace the capped $H$ by $K$.

The exact energy identity implies
\[
\energy{x_*-x_k}^2
\le\energy{x_*-y}^2+2|(b-Ax_k)^T(y-x_k)|.
\]
Since $A^2\preceq A$, \eqref{eq:exactenergy} and
\eqref{eq:crossbound} yield
\[
\norm{b-Ax_k}^2\le4K\vartheta^{2k}+\omega D^{12k+20}.
\]
The chosen $k_C$ makes the first term at most $\eps^2/4$:
use
\[
\vartheta\le\exp(-2/\sqrt{\bar K}),\qquad
k_C\ge\sqrt{\bar K}\log_2(16\bar K/\eps).
\]
The second term is at most $\eps^2/16$. Thus $\norm{b-Ax_k}<\eps$,
which implies \eqref{eq:eta} because its denominator is at least one.
This contradicts failure.

\subsection*{Case 3: the universal backward horizon $m=k_U$}
\label{sec:universal}
Put $k=m$. We use the following exact polynomial fact. For the even
integer $\ell\in\{k,k+1\}$, define
\begin{equation}\label{eq:reciprocalpoly}
G(t)=\frac{1-T_\ell(2t-1)}{2\ell^2},\qquad
\pi(t)=\frac{t-G(t)}{t^2}.
\end{equation}
The numerator has a double zero at zero; $\pi$ has degree at most
$\ell-2<k$, is positive on $[0,1]$, and satisfies
\begin{equation}\label{eq:recipbound}
\left|t-\frac1{\pi(t)}\right|\le\frac3{k^2-1},\qquad0\le t\le1.
\end{equation}
These are Lemmas 12--13 of Derezi\'nski, Nakatsukasa, and
Rebrova\cite{minberr}. This is the only external approximation lemma
used in this case. The shifted Chebyshev recurrence gives coefficient
sum at most $7^\ell$, so
\[
\coefnorm\pi\le16^{k+1}\le D^{k+1}.
\]
Also \eqref{eq:recipbound} gives $1/\pi(t)\le2$, so $\pi(t)\ge1/2$.
Therefore $y=\pi(A)b$ satisfies
\begin{equation}\label{eq:yberr}
\norm y\ge\tfrac12,\qquad
\berr(y):=\frac{\norm{b-Ay}}{\norm y}\le\frac3{k^2-1}.
\end{equation}
The spectral implication is immediate by applying
$A-\pi(A)^{-1}$ to $y$.

The failure assumption imposes an incompatible lower bound on every
vector in the computed-iterate span. Indeed, for
$z=\sum_{j=1}^k c_jx_j$,
\begin{equation}\label{eq:spanresidual}
b-Az=\left(1-\sum_{j=1}^k c_j\right)r_0
       +\sum_{j=1}^k c_jr_j+\sum_{j=1}^k c_jw_j.
\end{equation}
The Gram matrix of normalized $r_0,\ldots,r_k$ has off-diagonal entries
at most
\[
\frac{64E_k}{49\eps^2}\le\frac4{49}D^{-4k-12},
\]
so its off-diagonal row sums are below $1/4$.
In addition, \eqref{eq:Wbound} and \eqref{eq:failR} give
$W_j/R_j\le1/(8\sqrt k)$. Write
$Z=(\sum_{j=1}^k c_j^2R_j^2)^{1/2}$.
The first two terms in \eqref{eq:spanresidual} have norm at least
$\sqrt{3/4}\,Z$, while the last has norm at most $Z/8$.
Thus $\norm{b-Az}\ge Z/2$.
The failure bound $X_j/R_j\le2/\eps$ and Cauchy--Schwarz give
$\norm z\le2\sqrt k\,Z/\eps$. Every nonzero vector in this span
therefore satisfies
\begin{equation}\label{eq:spanberrlower}
\berr(z)\ge\frac{\eps}{4\sqrt k}.
\end{equation}

Lemma~\ref{lem:span} supplies a vector in the same span with
\[
\norm{z-y}\le\omega D^{3k+2}\le\frac{\eps}{128\sqrt k}.
\]
Relative to $\norm y\ge1/2$, this is at most
$t_0=\eps/(64\sqrt k)$, so $z\ne0$.
The definition of $k_U$ gives
\[
\frac3{k^2-1}\le\frac4{k^2}\le\frac{\eps}{16\sqrt k}.
\]
Perturbing \eqref{eq:yberr} now gives
\[
\berr(z)\le\frac{\berr(y)+t_0}{1-t_0}
\le\frac{\eps}{8\sqrt k},
\]
contradicting \eqref{eq:spanberrlower}.

\medskip
All possible minimizers of the horizon lead to a contradiction. Thus a
successful iterate is produced by $m$. Condition \eqref{eq:smallness}
is exactly implied by
\[
u\le\frac{\eps^2}{1024(n+1)K D^{12m+20}}.
\]
Every larger precision satisfies the same inequality, and the argument
applies to all real normalized inputs. This proves the full threshold
quantifier, \eqref{eq:mainupper}, and its order form
\eqref{eq:orderupper}.
\end{proof}

\section{Consequences with matching bounds}
\label{sec:consequences}

\begin{corollary}[Bounded conditioning]\label{cor:K2}
Uniformly for $n\ge1$, $1\le K\le2$, and $0<\eps<1/2$,
\[
\PCG(n,K,\eps)=\Theta\!\left(\log_2\frac n\eps\right).
\]
\end{corollary}
\begin{proof}
Here $H\le2$ and $D\le4096$ are universally bounded.
Moreover $m\le k_C=2\lceil\log_2(32/\eps)\rceil$.
Equation~\eqref{eq:mainupper} is therefore
$O(\log(n/\eps))$. The lower bound is \eqref{eq:loweridentity}.
\end{proof}


\begin{corollary}[Additive characterization at fixed tolerance]\label{cor:fixedtol}
For every fixed $0<\eps<1/2$, uniformly over all $n\ge2$ and $K\ge1$,
\[
\PCG(n,K,\eps)=\log_2(nK)+O_\eps(1).
\]
An explicit lower additive constant is $-5$. An upper additive constant
may be taken as
\[
C_\eps=1+\log_2\frac{1536}{\eps^2}
+(12k_U+20)\log_2\left(1024(64/\eps^2)^2\right).
\]
\end{corollary}
\begin{proof}
The lower inequality is \eqref{eq:jointlower}. For the upper inequality,
use $H\le64/\eps^2$, $m\le k_U$, and $n+1\le3n/2$ in
\eqref{eq:mainupper}, and allow one bit for its ceiling. Every term
remaining after $\log_2(nK)$ depends only on $\eps$.
\end{proof}

\begin{corollary}[Leading conditioning coefficient]\label{cor:largeK}
For every fixed $n\ge2$ and fixed $0<\eps<1/2$,
\[
\PCG(n,K,\eps)=\log_2K+O_{n,\eps}(1),\qquad
\lim_{K\to\infty}\frac{\PCG(n,K,\eps)}{\log_2K}=1.
\]
\end{corollary}
\begin{proof}
Once $K\ge64/\eps^2$, the quantity $H=64/\eps^2$ no longer depends
on $K$. Since $m\le n$, \eqref{eq:mainupper} is at most
$\log_2K+C_{n,\eps}$ for a finite constant independent of $K$.
On the other hand, \eqref{eq:lowerbreakdown} is at least
$\log_2K-1$. These two inequalities prove both assertions.
\end{proof}

This fixed-tolerance limit is compatible with the recovered
$n\log K$ lower family: that family decreases $\eps$ as $K$ increases.
Interchanging these parameter regimes would give an incorrect conclusion.

\begin{corollary}[Two dimensions and very fine tolerances]\label{cor:two}
Equation~\eqref{eq:two} holds uniformly. More generally, if $n\ge2$ and
$\log_2(1/\eps)\ge n\log_2(2K)$, then
\[
\PCG(n,K,\eps)=\Theta\!\left(\log_2\frac{nK}{\eps}\right)
\]
with universal constants on this region.
\end{corollary}
\begin{proof}
Use $H\le K$ and $m\le n$ in \eqref{eq:mainupper}.
When $n=2$, its logarithm is $O(\log K+\log(1/\eps))$.
Under the second displayed condition, the term $n\log(2K)$ is
absorbed by $\log(1/\eps)$. Theorem~\ref{thm:lower} supplies the
matching lower bounds.
\end{proof}

At $K=1$, the tighter first-step estimate from the recovered manuscript
is $\PCG\le\lceil\log_2(128(n+1)/\eps)\rceil$.
Together with \eqref{eq:loweridentity}, this yields the additive form
$\PCG(n,1,\eps)=\log_2(n/\eps)+O(1)$ with universal additive constants.
The exact scalar result is stronger still when $n=1$.

\section{The retained hard family and the unresolved gap}
\label{sec:gap}

The recovered manuscript proves the following separate lower-bound
family. Let $s=n-1$, and let $M_s$ have diagonal $(1,2,\ldots,2)$ and
adjacent off-diagonal entries $-1$. Its inverse bidiagonal factor gives
$\lambda_{\min}(M_s)\ge1/s^2$, and $\norm{M_s}<4$.
Take
\[
T=2^\ell,\qquad A=M_s\oplus[T],\qquad b=(e_1,1)^T,
\qquad p=s\ell.
\]
The only nonzero operation error is $+2^{-p}$ on the last addition in
$\rho_0$. All later operations are exact. In particular the true and
recursive residuals agree. For each fixed $n\ge2$, the recovered
recurrence analysis shows that, for all sufficiently large integers
$\ell$,
\begin{equation}\label{eq:retainedhard}
\PCG\!\left(n,s^2 2^\ell,\frac1{64s^2 2^\ell}\right)>s\ell.
\end{equation}
The quantifier is fixed $n$, followed by sufficiently large $\ell$;
no dimension-uniform threshold is supplied. The full proof remains in
Section 5 of the preserved prior manuscript. Its finite exact witnesses
have been rerun in this package. The new results do not strengthen the
quantifier in \eqref{eq:retainedhard}.

\paragraph{A concrete gap still present.}
Consider $n\ge2$ with
\begin{equation}\label{eq:gaptriple}
K=n^2,\qquad\eps=n^{-2}.
\end{equation}
The generally applicable lower bound established in this package is
$\Omega(\log n)$, while its upper bound is $O(n\log n)$.
The retained large-outlier limit does not close this growing-dimension
regime. Neither an $\Omega(n\log n)$ lower bound nor an
$O(\log n)$ upper bound on \eqref{eq:gaptriple} is proved here.

A full solution must supply a universally matching characterization on
such regimes as well as on the sharp regions already covered. The
upper bound in this manuscript is not asserted to be sharp everywhere.
In particular, its reciprocal-polynomial horizon is a sufficient
precision budget, not a matching lower bound when $n$ exceeds that
horizon. An interval-polynomial obstruction is also not, by itself, an
admissible CG execution; the actual recurrence and its spectral weights
must be checked.

The reason for the partial label is this specific mathematical gap,
not merely the problem's status in a repository. No full-resolution or
priority claim is made.

\section{Exact checks, reproducibility, and proof scope}
\label{sec:checks}

All executable checks use Python's \texttt{fractions.Fraction} and the
standard library. No approximate square root or floating logarithm is
used to decide a certificate. The error-envelope executor checks input
significands, symmetry, exact $LDL^T$ positivity, and every supplied
operation-error magnitude. It includes the required zero terms and
separates products from additions.

\subsection{What was checked}

The regenerated reports record:
\begin{center}
\begin{tabular}{p{0.70\textwidth}r}
\toprule
Check & Count\\
\midrule
Scalar first-step sign-corner executions (three precisions) & 1,536\\
First-step executions at a certified precision (finite inputs/schedules) & 256\\
Exact reciprocal-polynomial constructions & 31\\
Rational-grid reciprocal-polynomial inequalities & 3,999\\
Parameter-triple consistency checks for bound evaluation & 90\\
Joint dimension/conditioning breakdown witnesses & 14\\
Dense matvec prefix-bound executions & 24\\
Retained exact counterexample witnesses & 15\\
Unit tests, including the retained regression suite & 26\\
\bottomrule
\end{tabular}
\end{center}
The first-step executions use diagonal SPD inputs with known exact
spectral norm, dimensions $2,3,5,8$, condition bounds
$1,2,2^6,2^{20}$, two coarse tolerances, and eight deterministic error
schedules. These are finite checks, not an exhaustive multivariate
error search. The scalar corner enumeration is exhaustive only for its
nine corner signs; its continuous-box validity is proved analytically
in Theorem~\ref{thm:scalar}.

For known rational $a=\norm A$ and squared norms $R^2,X^2,B^2$, testing
$\eta\le\eps$ reduces to comparing
\[
C=R^2/\eps^2-a^2X^2-B^2
\quad\text{with}\quad 2a\sqrt{X^2B^2}.
\]
If $C<0$, success is strict. If $C\ge0$, compare the squares, retaining
the zero case separately. This gives exact comparison without extracting
a square root. The older hard-family verifier uses an independent
sufficient failure certificate based on
$(aX+B)^2\le2(a^2X^2+B^2)$.

\subsection{Commands and bound evaluation}

From the package root, run
\begin{verbatim}
python -m unittest discover -s code -v
python code/verify_extended.py
python code/verify_legacy.py --output results/legacy_checks.json
python code/bounds.py --n 16 --K 256 --epsilon 1/256
\end{verbatim}
Run without \texttt{-O}; the witness scripts reject optimized mode
because they use assertions. Python 3.10 or later is sufficient.

The scalar evaluator returns the exact integer threshold. The general
evaluator reports separate lower and upper bounds, not an alleged exact
answer. To avoid constructing enormous rational powers, its conservative
integer version of \eqref{eq:mainupper} is
\[
\left\lceil\log_2\frac{1024(n+1)K}{\eps^2}\right\rceil
+(12m+20)\lceil\log_2 D\rceil.
\]
This can exceed the ceiling in the theorem by at most $12m+20$ bits.
It remains a valid sufficient bound and has the same stated order.
The one-step evaluator compares the square of
$\eps-512(n+1)K2^{-p}$ with $M(K)^2$, after checking the sign.

\subsection{What the checks do not establish}

Passing these checks does not prove the bootstrap for arbitrary inputs,
the uniform reciprocal-polynomial inequality between grid points, the
large-outlier asymptotics in arbitrary dimension, or a matching precision
formula. The bootstrap, transfer arguments, and corollaries are written
mathematical proofs; the reciprocal-polynomial inequality has a cited
analytic proof. None is represented as formally verified or independently
reviewed. The preserved old archive is also explicitly partial.

\appendix
\section{A compact proof-dependency audit}

The scalar theorem depends only on counting the specified nine operations
and monotonicity of a rational expression. The first-step theorem depends
on the moment inequality and local scalar envelopes. The main upper bound
uses the following chain, with no circular assumption of successful
termination:
\begin{center}
\begin{tabular}{p{0.26\textwidth}p{0.65\textwidth}}
\toprule
Stage & Required fact\\
\midrule
Local equations & Actual $K$ controls denominator perturbations.\\
Current step legality & Earlier defect bounds give $r_j^Tv_j>0$.\\
Capped step length & Failure of the newly produced iterate rules out a long step.\\
Bootstrap propagation & Coefficient bounds, local errors, and adjacent cancellations.\\
Horizon contradiction & Rank, Chebyshev energy comparison, or reciprocal-polynomial span comparison.\\
Threshold conclusion & The same sufficient inequality holds at every larger precision.\\
\bottomrule
\end{tabular}
\end{center}
The joint breakdown lower bound is separate from this chain: it uses
exactly representable inputs and rational, greedy accumulation factors.

The only non-elementary external approximation fact used in the new upper
proof is \eqref{eq:recipbound}. Removing its branch leaves the still-valid
upper bound with $m=\min\{n,k_C\}$, including the cap, the bounded-$K$
order, and the leading large-$K$ coefficient.

\begin{thebibliography}{9}
\bibitem{ie20}
OpenProblemsInNLA repository.
\emph{IE-20: Precision required for conjugate gradients to attain backward
accuracy in n steps}. Specification accessed September 13, 2026.
\url{https://github.com/ajt60gaibb/OpenProblemsInNLA/tree/main/linear-systems-and-elimination/IE-20}.
Readable source:
\url{https://raw.githubusercontent.com/ajt60gaibb/OpenProblemsInNLA/main/linear-systems-and-elimination/IE-20/README.md}.

\bibitem{minberr}
M. Derezi\'nski, Y. Nakatsukasa, and E. Rebrova.
\emph{Towards Universal Convergence of Backward Error in Linear System Solvers}.
arXiv:2604.16075v2, May 22, 2026. Lemmas 12--13 and Theorem 9.
\url{https://arxiv.org/html/2604.16075v2}.

\bibitem{bcm}
T. Bake, E. Carson, and Y. Ma.
\emph{Forward and backward error bounds for a mixed precision
preconditioned conjugate gradient algorithm}.
arXiv:2510.11379v3, July 21, 2026. Lemma 6 and Theorem 2.
\url{https://arxiv.org/html/2510.11379v3}.

\bibitem{sinha}
M. Sinha.
\emph{Finite-Precision Symmetric Krylov Methods: Exact Rounding Examples,
Block Paige Identities, and a Variable-Block Lanczos Model}.
arXiv:2609.08066v1, September 8, 2026. Section 5.
\url{https://arxiv.org/html/2609.08066v1}.
\end{thebibliography}
\end{document}
